\chapter{Introduction}
\label{chap:intro}

Quantum computing matters because it addresses computational tasks for which classical methods appear fundamentally inadequate. Feynman famously argued that the state space of a generic quantum many-body system grows exponentially with system size, so that classical computers cannot efficiently simulate quantum physics in full generality, whereas a controllable quantum device should be able to do so more naturally~\cite{Feynman1982,Lloyd1996}. This prospect makes quantum computing especially compelling for quantum simulation, with anticipated applications ranging from molecular electronic structure, ground-state energies, and reaction dynamics relevant to chemistry and drug design~\cite{Cao2019,McArdle2020}, to strongly correlated materials, high-temperature superconductivity, and lattice gauge theories in physics~\cite{Georgescu2014,Wiese2013}, and to emerging problems in computational biology and medicine such as protein folding, photosynthetic energy-transfer models, and quantum-assisted drug discovery~\cite{Outeiral2021,Santagati2024}. Beyond simulation, quantum algorithms can also provide rigorous advantages for certain classical problems. Shor's algorithm yields an exponential speedup for integer factorization and discrete logarithms, thereby threatening widely used public-key cryptosystems and motivating the ongoing transition to post-quantum cryptography~\cite{Shor1994,Dalzell_2025,NIST2024PQC}, while Grover's algorithm offers a quadratic query speedup for unstructured search~\cite{Grover1996,Montanaro2016}. Recent surveys further identify finance as an important application area for fault-tolerant quantum algorithms, with proposed uses in portfolio optimization, Monte Carlo derivative pricing, and risk analysis~\cite{Dalzell_2025,Orus2019,Herman2023}. Other directions, including quantum machine learning and quantum optimization, further broaden the field's scope, although their practical advantage remains an active research question~\cite{Biamonte2017,Blekos2024}. Together, these scientific and computational prospects have driven strong worldwide efforts, as well as substantial public and private investment, to build practical quantum computing platforms~\cite{NaturePhysics2023,Soller2025}.

A practical quantum computer must be realized in a physical system of effective two-level degrees of freedom, or qubits, that can be initialized, coherently controlled, entangled, and measured with high fidelity, while maintaining coherence for times long compared with the duration of gate operations \cite{DiVincenzo2000,Ladd2010}. Several hardware platforms satisfy these requirements to different degrees. Nitrogen-vacancy centers in diamond offer exceptionally long coherence times even at room temperature, but scalable fabrication and controlled interconnects remain challenging \cite{Doherty2013,Pezzagna2021}. Trapped ions provide some of the best coherence times and gate fidelities demonstrated so far, although gate operations are relatively slow and large-scale integration remains technically demanding \cite{Monroe2013,Bruzewicz2019}. Neutral-atom processors have demonstrated large qubit arrays together with dynamically reconfigurable connectivity, making them attractive for both programmable quantum simulation and gate-based computation \cite{Saffman2016,Bluvstein2022}. Photonic platforms are especially natural for quantum communication and networking, but deterministic two-qubit gates remain difficult because photons interact only weakly in ordinary optical media \cite{Kok2007,Slussarenko2019}. Superconducting circuits, by contrast, combine fast gate operations with lithographically scalable fabrication and mature microwave control techniques \cite{Krantz2019,Kjaergaard2020}. No single platform has yet simultaneously optimized all of the requirements for large-scale quantum computation, and current hardware development is therefore shaped by trade-offs among coherence, connectivity, control speed, and manufacturability \cite{DiVincenzo2000,Ladd2010}. Among gate-based implementations, superconducting circuits are currently one of the most mature platforms, as reflected both in major public roadmaps from IBM and Google and in recent large-scale error-correction demonstrations on superconducting processors \cite{IBMRoadmap2025,GoogleRoadmap2026,GoogleQEC2025,Kjaergaard2020}. This thesis therefore focuses exclusively on superconducting circuits.

Superconducting qubits encode quantum information in quantized macroscopic variables of Josephson circuits, where nonlinear circuit elements allow selected levels of an electrical mode to be isolated and controlled as a qubit \cite{Krantz2019,Kjaergaard2020}. Over the past two decades, a wide variety of designs has emerged, each optimizing a different balance between anharmonicity, tunability, coherence, and noise protection. The charge qubit, or Cooper-pair box, provided the first clear demonstration of coherent control in a superconducting qubit, but its strong sensitivity to offset-charge fluctuations limited its coherence \cite{Nakamura1999,Pashkin2009}. The flux qubit instead encodes information in superpositions of opposite persistent-current states, offering strong anharmonicity and convenient magnetic control, although flux noise becomes a central decoherence channel \cite{Mooij1999,Yan2016}. The transmon, obtained by shunting the Cooper-pair box with a large capacitor, exponentially suppresses charge noise at the cost of reduced anharmonicity \cite{Koch2007,Houck2009}. Fluxonium and especially heavy fluxonium recover large anharmonicity while supporting very long coherence through the use of a superinductor \cite{Manucharyan2009,Zhang2021}. Related protected designs, including the bifluxon and the 0-$\pi$ qubit, seek hardware-level suppression of dominant noise channels by engineering symmetry and disjoint wave-function support into the circuit itself \cite{Kalashnikov2020,Gyenis2021}. More exotic proposals, such as current-mirror and \(\cos 2\phi\)-based protected circuits, extend this strategy further \cite{Kitaev2006,Smith2020}. Despite this zoology, the transmon has become the workhorse of both academia and industry because it combines conceptual simplicity, reliable fabrication, high-fidelity control, and fast microwave readout in a scalable architecture \cite{Kjaergaard2020,Jiang2025}. In present-day devices, this platform supports coherence times from tens to hundreds of microseconds together with high-fidelity gate operations, and it has been scaled from 105-qubit-class processors such as Google's Willow chip to IBM's 1,121-qubit Condor processor \cite{GoogleQEC2025,IBMCondor2023}. This thesis therefore focuses on superconducting circuits, with particular emphasis on the transmon architecture.

Alongside transmon-based processors, superconducting microwave cavities provide a distinct and complementary paradigm for quantum information processing. In both three-dimensional and planar implementations, a cavity is a linear superconducting resonator that can store microwave photons with coherence times far exceeding those of typical transmon qubits, while universal control is usually supplied by coupling the cavity to a nonlinear ancilla such as a transmon \cite{Krasnok2024,Copetudo2024}. Early progress in three-dimensional circuit QED established the exceptional coherence of machined cavities \cite{Paik2011,Reagor2013}, and continued improvements in materials, fabrication, and surface treatment pushed 3D cavity photon lifetimes from the millisecond regime to tens of milliseconds and ultimately beyond \(2\,\mathrm{s}\) in state-of-the-art niobium resonators at millikelvin temperatures \cite{Reagor2016,Milul2023,Romanenko2020}. Planar two-dimensional superconducting resonators generally suffer stronger dielectric and surface-interface losses and therefore do not yet match the best 3D cavities \cite{Wang2009,Megrant2012}; nevertheless, recent on-chip quantum-memory designs have exceeded the millisecond scale, demonstrating that 2D architectures can also reach remarkably long coherence when carefully engineered \cite{Ganjam2024,Krayzman2024}. The underlying physical advantage of both 2D and 3D cavities is that, as nearly ideal linear oscillators, they avoid the dominant charge, flux, and critical-current noise channels associated with Josephson qubits, so photon loss is typically the primary error mechanism while pure dephasing remains comparatively weak \cite{Joshi2021,Cai2021}. This strong bias toward photon-loss errors, together with the large Hilbert space of a harmonic oscillator, makes superconducting cavities especially attractive for bosonic encoding and quantum error correction.

A superconducting cavity by itself is only a harmonic oscillator, with equally spaced energy levels, so linear drives cannot selectively address individual transitions or generate the non-Gaussian operations required for universal quantum control \cite{Blais2021,Ma2021}. In circuit QED, the standard solution is to couple the cavity dispersively to an anharmonic ancilla, typically a transmon, so that the qubit and cavity frequencies acquire state-dependent shifts characterized by the dispersive coupling \(\chi\) \cite{Blais2021,Schuster2007}. This interaction enables a rich toolbox for bosonic control. In the strong-dispersive regime, the selective number-dependent arbitrary phase (SNAP) gate applies independently programmable phases to chosen Fock states by driving ancilla transitions at frequencies \(\omega_q+n\chi\); together with cavity displacements, SNAP gates provide universal single-cavity control \cite{Heeres2015,Krastanov2015}. Their reliance on spectrally selective ancilla pulses, however, makes them increasingly vulnerable to Stark shifts and spectator-mode crosstalk in multimode devices \cite{You2024,Ma2021}. Echoed conditional displacement (ECD) gates address the same problem differently: by combining conditional cavity displacements with ancilla echo pulses, they realize fast universal control even when the bare dispersive coupling is weak, and because the ancilla operation is not photon-number selective they are generally less sensitive than SNAP to crosstalk, although not completely immune to it \cite{Eickbusch2022,You2024}. More recently, conditional-NOT displacement (CNOD) has extended this philosophy to multimode architectures, showing that a single ancilla can efficiently control and entangle several weakly coupled oscillators \cite{Diringer2024,Copetudo2024}. In parallel, optimized SNAP-displacement sequences and improved cavity-control protocols have further reduced coherent errors and expanded the accessible family of bosonic states \cite{Kudra2022,Landgraf2024}. For two-cavity operations, however, one generally needs an activated bilinear coupling between modes. Early demonstrations realized programmable beam-splitter interactions through a driven transmon ancilla \cite{Gao2018,Zhang2019}, whereas newer architectures employ dedicated parametric couplers such as SNAIL-based three-wave mixers and parity-protected converters to obtain fast swaps, high on/off ratios, and low idle-time backaction on the memories \cite{Chapman2023,Lu2023}. More recently, linear quantum couplers such as the LINC have been proposed and realized to push this idea further by remaining effectively linear when idle while activating cleaner parametric interactions under drive \cite{Maiti2025,Pietikainen2024}. Together, these developments define the modern control toolbox for bosonic circuit QED and have enabled beam-splitter-based logical operations, including dual-rail cavity encodings and error-detectable entangling gates between bosonic qubits \cite{Teoh2023,Tsunoda2023}.

Superconducting cavities are especially attractive for bosonic quantum error correction because their large Hilbert space and strongly biased noise channels allow a single oscillator mode to store a logical qubit with hardware-efficient redundancy \cite{Cai2021,Copetudo2024}. The central idea is to encode logical information in nonclassical superpositions of Fock states or phase-space states chosen so that the dominant physical error---typically single-photon loss---maps the system into a distinguishable error manifold that can be detected and corrected before the logical information is destroyed \cite{Cai2021,Copetudo2024}. Among the major bosonic code families, cat codes encode information in superpositions of coherent states \(\ket{\pm \alpha}\), for which single-photon loss changes the photon-number parity and appears as a logical bit-flip-type error \cite{Leghtas2013,Mirrahimi2014}. Binomial codes instead use carefully engineered finite superpositions of Fock states and can be designed to correct photon loss, photon gain, and dephasing; the lowest-order example is \(\ket{0_L} = (\ket{0}+\ket{4})/\sqrt{2}\) and \(\ket{1_L} = \ket{2}\) \cite{Michael2016,Albert2018}. Gottesman--Kitaev--Preskill (GKP) codes encode a logical qubit in grid states of oscillator phase space and are tailored to correct small displacement errors in both quadratures \cite{Gottesman2001,Brady2024}. A complementary dual-rail encoding stores the logical qubit nonlocally across two cavities, with \(\ket{0_L}=\ket{01}\) and \(\ket{1_L}=\ket{10}\); in this case a photon-loss event maps the state outside the code space as an erasure-type error that is comparatively easy to detect \cite{Teoh2023,Copetudo2024}. Experimentally, bosonic QEC in superconducting cavities has progressed from the first break-even cat-code experiment \cite{Ofek2016,Cai2021}, through the first full GKP error-correction demonstrations in circuit QED \cite{CampagneIbarcq2020,Brady2024}, to beyond-break-even real-time correction with GKP states \cite{Sivak2023,Brady2024} and with binomial codes \cite{Ni2023,Cai2021}. Cat-code experiments have likewise demonstrated extended logical lifetimes and exponentially suppressed bit-flip rates as the phase-space separation is increased \cite{Gertler2021,Lescanne2020}. This progress relies critically on noise bias: bosonic encodings are most powerful when photon loss remains the dominant cavity error, because then increasing the encoding size can strongly reduce logical failure rates, whereas comparable dephasing substantially weakens that advantage \cite{Albert2018,Leviant2022}. For that reason, suppressing cavity dephasing---even at the cost of additional hardware or control complexity---remains central to the bosonic-QEC program \cite{Cai2021,Leviant2022}.

One of the prominent techniques for implementing quantum optimal control is Gradient Ascent Pulse Engineering (GRAPE).%
\nomenclature{GRAPE}{Gradient Ascent Pulse Engineering}
The central goal of optimal control is to adjust control parameters in such a way that the infidelity of the desired state transfer or gate operation is minimized. In GRAPE, the minimization is based on gradient descent and thus generally requires the evaluation of gradients. Automatic differentiation (AD)~\cite{baydin2018automatic} is a convenient tool that has also made an impact on quantum optimal control with GRAPE~\cite{pacleung2017speedup,pacqoc,abdelhafez2019gradient,grapead}. Internally, AD builds a computational graph and applies the chain rule to automatically compute the gradient of a given function.

However, the convenience of AD comes at the cost of memory required for storing the full computational graph. The use of semi-automatic differentiation (semi-AD)~\cite{semiadgoerz2022quantum} partially reduces the memory overhead compared to full automatic differentiation (full-AD). This reduction can still be insufficient to tackle optimal control for quantum systems of rapidly growing size~\cite{arute2019quantum,gong2021quantum}. This motivates a revisit of the original GRAPE scheme~\cite{nmrkhaneja2005optimal}, which is based on hard-coded gradients (HG) and has the smallest possible memory cost. In this thesis, I perform a scaling analysis of memory usage and runtime and develop a decision tree guiding the optimal choice of strategy among HG, semi-AD, and full-AD. I exemplify the scaling behavior with benchmarks for concrete optimization tasks.

At the application level, quantum optimal control helps build up the quantum processor unit (QPU).%
\nomenclature{QPU}{Quantum Processor Unit}
In the existing superconducting QPU architecture, the transmon can serve as either a qubit or an ancilla for a bosonic qubit. In both cases, Rabi oscillation in the first two-level subspace of the transmon is crucial for realizing universal gates. Conventionally, to realize such oscillation in a bare transmon (without coupling to other systems), a Gaussian pulse is applied where the drive frequency is chosen to be on resonance with the transmon frequency. In a QPU, the transmon is dispersively coupled with other systems, causing a Stark-shifted frequency that is state dependent. In such a case, intuitively, the drive frequency can be optimized to adapt to this frequency shift. However, such state-dependent drive frequency can cause large hardware overhead. Another more feasible choice is to fix the drive frequency and optimize the pulse to improve the robustness of the Rabi oscillation fidelity against detuning frequency. Such robust pulses are hardly obtained analytically, and quantum optimal control can be employed to find them.

Based on the analysis of the numerical implementation of GRAPE, I develop a resource-efficient GRAPE-based optimization procedure tailored to generate detuning-robust pulses. I benchmark the optimized pulses against conventionally used Gaussian pulses, demonstrating improved robustness of transmon gate fidelity.

\section{Thesis outline}

This thesis is organized as follows.
Chapter~\ref{chap:grape} introduces the GRAPE framework and derives hard-coded analytical gradients for key cost contributions used in state-transfer and gate-optimization problems. The chapter also discusses the numerical evaluation of these gradients using the scaling-and-squaring method and the auxiliary matrix method, culminating in a memory-efficient forward-backward propagation scheme.

Chapter~\ref{chap:scaling} presents a scaling analysis of runtime and memory usage for hard-coded gradients, full-AD, and semi-AD. Benchmark studies for state-transfer and gate-optimization problems illustrate the scaling behavior. A decision tree is developed to guide the choice of the optimal numerical strategy.

Chapter~\ref{chap:robust} turns to an application of GRAPE: the design of detuning-robust control pulses for transmon gates. The optimization procedure is detailed, and the resulting pulses are benchmarked against conventional DRAG pulses and pulses from Q-CTRL.

Chapter~\ref{chap:conclusion} summarizes the key findings and discusses opportunities for future research.

Appendix~\ref{app:muscaling} discusses the scaling of the number of matrix-vector multiplications $\mu$ with Hilbert space dimension for several concrete Hamiltonians. Appendix~\ref{app:matdiag} presents an alternative hard-coded gradient scheme based on matrix diagonalization and its scaling analysis.
